\section{Recommended Future Work}
\label{sec:future_work}

The system is functional and the headline numbers are consistent, but there
are seven directions in which the work would extend naturally. They are
ordered roughly by ratio of expected impact to effort.

\subsection{Human Evaluation}
The strongest single follow-up is a small-N human evaluation by trained MI
clinicians, even on the order of 20--30 dialogues, to calibrate the
LLM-as-judge scores to clinician judgement. The current evaluation is a
proxy; the next claim worth making about this system is one that survives
contact with a human rater.

\subsection{Larger Base Model}
The Gemma-2-2B base was chosen for local CPU deployability, and the
evaluation shows a substantial residual gap to the Claude Haiku reference.
A direct experiment is to fine-tune Gemma-2-9B (or a similarly sized
\texttt{Llama-3.1-8B-Instruct}) with the same LoRA recipe on the same data
and measure how much of the remaining gap closes. The trade-off is latency:
9B-class models are GPU-required for interactive use, so this becomes a
hosted-inference deployment rather than a local-CPU one.

\subsection{Embedding-Based Safety Classifier (as a Layer, not a Replacement)}
The current safety screen is regex-based by deliberate choice for
auditability. But regex coverage is enumerative -- semantically equivalent
phrasings can slip through. A reasonable architecture is to add an
embedding-based classifier as a \emph{second} safety layer, run only when
the regex screen passes; if the classifier flags the turn at high confidence,
the system would conservatively short-circuit to the crisis response. This
preserves regex auditability for the explicitly-listed cases while widening
coverage for paraphrases. Training data is the bottleneck: a labeled crisis
corpus would be needed and is not freely available.

\subsection{Multi-Turn Coherence}
The current evaluation scores responses turn-by-turn against per-turn
references. MI is a multi-turn process -- summaries, change-talk
elicitation, and stage-of-change progression all unfold across many turns.
A natural extension is to evaluate at the conversation level: does the
system summarize at the right moments, does it elicit and then reflect change
talk, does it avoid bringing up barriers prematurely? This requires a
conversation-level rubric and a different judge prompt structure.

\subsection{Cross-Session Memory}
The system is currently stateless across sessions -- each conversation starts
fresh. For real use this is a meaningful gap, since MI is most effective
across multiple sessions where the counselor remembers prior commitments and
prior reasons-for-change. A practical next step is per-user persistent memory
of (a) explicitly stated reasons-for-change and (b) self-reported milestones,
gated by user consent and excluded from the model's general context window
to avoid prompt bloat.

\subsection{Latency}
The 2--4~s per-turn latency on CPU is workable for demos but breaks the
conversational flow of MI in real-time use. Three independent tactics would
reduce this: (a) further quantization (Q4\_0 or Q3\_K\_S) at a small quality
cost, (b) speculative decoding via a smaller draft model, and (c) GPU
deployment with continuous batching. (a) is the cheapest experiment.

\subsection{Broader Domain Coverage}
The project is scoped to alcohol use specifically. The same pipeline -- swap
the knowledge base, refit the synthetic generator's scenario types, and
re-fine-tune -- should generalize to other behavioral-health domains for
which MI is appropriate (smoking cessation, gambling, dietary change). This
is mostly an integration question rather than a research one. The interesting
research question alongside it is how much of the alcohol-trained model's
behavior transfers across domains zero-shot, and whether a single
multi-domain fine-tune outperforms domain-specific ones at the size budget
considered here.
